\documentclass[../../../include/open-logic-section]{subfiles}

\begin{document}

\olfileid{sth}{ordinals}{opps}
\olsection{后继序数与极限序数}

在接下来的几章中，我们将大量使用序数。因此，先引入一些简单的概念会很有帮助。

\begin{defn}
对于任意序数 $\alpha$，其\emph{后继}是 $\ordsucc{\alpha}
=\alpha \cup \{\alpha\}$。若存在序数~$\beta$ 使得 $\ordsucc{\beta} = \alpha$，则称 $\alpha$ 为\emph{后继}序数。称 $\alpha$ 为\emph{极限}序数，当且仅当 $\alpha$ 既不是空集也不是后继序数。
\end{defn}

\noindent
以下结果表明，这是\emph{后继}的正确概念：

\begin{prop}
对任意序数 $\alpha$：
\begin{enumerate}
	\item $\alpha \in \ordsucc{\alpha}$；
	\item $\ordsucc{\alpha}$ 是一个序数；
	\item 不存在序数 $\beta$ 使得 $\alpha \in \beta \in \ordsucc{\alpha}$。
	\end{enumerate}
\end{prop}

\begin{proof}
显然，$\alpha \in \alpha \cup \{\alpha\} = \ordsucc{\alpha}$。
此外，$\ordsucc{\alpha}$ 是由序数组成的传递集，因此由 \olref[basic]{corordtransitiveord} 知它是一个序数。
且不可能 $\alpha \in \beta \in \ordsucc{\alpha}$，因为若不然，要么 $\beta \in \alpha$ 要么 $\beta = \alpha$，这与 \olref[basic]{ordordered} 矛盾。
\end{proof}

\noindent
这也使我们可以使用超限归纳法的一种变体：

\begin{thm}[简单超限归纳法]\ollabel{simpletransrecursion}
设 $\phi(x)$ 为一个满足以下条件的公式：
\begin{enumerate}
	\item $\phi(\emptyset)$；且
	\item 对任意序数 $\alpha$，若 $\phi(\alpha)$ 则 $\phi(\ordsucc{\alpha})$；且
	\item 若 $\alpha$ 是极限序数且 $(\forall \beta \in \alpha)\phi(\beta)$，则 $\phi(\alpha)$。
\end{enumerate}
那么 $\forall \alpha \phi(\alpha)$。
\end{thm}

\begin{proof}
我们证明其逆否命题。假设存在某个序数不满足 $\phi$；令 $\gamma$ 为这样的最小序数。那么要么 $\gamma = \emptyset$，要么存在满足 $\phi(\alpha)$ 的 $\alpha$ 使得 $\gamma = \ordsucc{\alpha}$；要么 $\gamma$ 是极限序数且 $(\forall \beta \in \gamma)\phi(\beta)$。
\end{proof}

\noindent
最后引入一个稍后将用到的记号：

\begin{defn}\ollabel{defsupstrict}
若 $X$ 是序数的集合，则 $\supstrict(X) = \bigcup_{\alpha
\in X} \ordsucc{\alpha}$。
\end{defn}

\noindent
这里，“lsub”代表“最小严格上界”。\footnote{有些书用“$\text{sup}(X)$”表示这个。但另一些书用“$\text{sup}(X)$”表示最小的\emph{非严格}上界，即 $\bigcup X$。如果 $X$ 有最大元 $\alpha$，这两个概念就不同了：最小的\emph{严格}上界是 $\ordsucc{\alpha}$，而最小的\emph{非严格}上界就是 $\alpha$。} 下面的结果解释了这一点：

\begin{prop}
若 $X$ 是序数的集合，则 $\supstrict(X)$ 是大于 $X$ 中每个序数的最小序数。
\end{prop}

\begin{proof}
令 $Y = \Setabs{\ordsucc{\alpha}}{\alpha \in X}$，于是 $\supstrict(X) = \bigcup Y$。由于序数是传递的，且序数的每个元素都是序数，故 $\supstrict(X)$ 是由序数组成的传递集，因此由 \olref[basic]{corordtransitiveord} 知它是一个序数。

若 $\alpha \in X$，则 $\ordsucc{\alpha} \in Y$，于是 $\ordsucc{\alpha}
\subseteq \bigcup Y = \supstrict(X)$，因此 $\alpha \in \supstrict(X)$。所以 $\supstrict(X)$ 严格大于 $X$ 中的每个序数。

反之，若 $\alpha \in \supstrict(X)$，则对某个 $\beta \in X$ 有 $\alpha \in \ordsucc{\beta} \in Y$，于是 $\alpha \leq \beta \in X$。所以 $\supstrict(X)$ 是 $X$ 的\emph{最小}严格上界。
\end{proof}

\end{document}